% =====================================================================
% FUNDAMENTAÇÃO TEÓRICA (Referencial Teórico)
% Reúna os conceitos necessários para o leitor compreender o trabalho.
% Este capítulo demonstra: citação (indireta, direta curta com \enquote, direta
% longa com o ambiente citacao, grifo nosso, supressão [...] e apud), uso correto
% de siglas (declarar uma vez, depois \gls{}), nota de rodapé, figura, equação,
% alíneas (ABNT) e listas.
% =====================================================================

Lorem ipsum dolor sit amet, consectetur adipiscing elit, sed do eiusmod tempor
incididunt ut labore et dolore magna aliqua~\parencite{exemplo-livro}. Ut enim ad
minim veniam, quis nostrud exercitation ullamco laboris.

% =====================================================================
% SIGLAS — A REGRA: DECLARE UMA VEZ, DEPOIS SEMPRE \gls{}
% =====================================================================
% NBR 14724:2024, 5.6: "A sigla, quando mencionada pela primeira vez no texto,
% deve ser indicada entre parênteses, precedida pelo nome completo" — ex.:
% Associação Brasileira de Normas Técnicas (ABNT). Nas menções seguintes, só a
% sigla.
%
% 1) DECLARE a sigla uma única vez, de uma destas duas formas (escolha uma):
%      a) no preâmbulo do thesis.tex:  \nomenclature[S]{ABNT}{Associação...}
%      b) aqui, na primeira menção:    \sigla{IBGE}{Instituto Brasileiro...}
%
% 2) A PARTIR DAÍ, toda menção é \gls{ABNT} / \gls{IBGE} — inclusive a primeira,
%    se você declarou pelo preâmbulo. O pacote cuida da regra 5.6 sozinho: na
%    primeira vez que a sigla aparece no texto ele imprime "nome completo
%    (SIGLA)"; da segunda em diante, apenas "SIGLA". Você NÃO precisa controlar
%    qual menção é a primeira, nem escrever o nome completo à mão.
%
% 3) NÃO digite a sigla como texto puro. Ela sai igual no PDF, mas perde o link
%    para a entrada na lista de abreviaturas e siglas — e, se for a primeira
%    menção, não expande (a regra 5.6 fica por sua conta).
%
% A chave é a própria sigla, então \gls{IBGE} para a sigla IBGE. Variações:
% \sigla*{}{} apenas cadastra, sem escrever nada no texto (útil para siglas que
% só aparecem em figuras ou tabelas), e \sigla[A]{}{} manda a entrada para o
% grupo "Abreviaturas" em vez de "Siglas".
As normas de apresentação de trabalhos acadêmicos são publicadas pela
\gls{ABNT} --- sigla declarada no preâmbulo, cuja primeira menção sai por
extenso automaticamente. Já os dados populacionais provêm do
\sigla{IBGE}{Instituto Brasileiro de Geografia e Estatística}, declarado na
própria linha em que é citado. Nas menções seguintes, ambas aparecem apenas
como sigla: \gls{ABNT} e \gls{IBGE}, cada uma com link para a entrada na lista.

%-------------------------------------------------------------
\section{Primeiro conceito}
\label{sec:conceito-um}
%-------------------------------------------------------------

Duis aute irure dolor in reprehenderit in voluptate velit esse cillum dolore eu
fugiat nulla pariatur~\parencite{exemplo-artigo}. A \autoref{fig:exemplo} ilustra
o conceito.

% Citação DIRETA curta (até 3 linhas): fica no corpo do texto, entre aspas, com a
% fonte e a página logo após (NBR 10520:2023). Use \enquote{...} (csquotes) para as
% aspas — NUNCA a aspa reta ", que sob XeLaTeX sai como ”...” (fecha nas duas pontas).
Nesse sentido, \enquote{a citação direta de até três linhas permanece no corpo do
texto, entre aspas}~\parencite[p. 15]{exemplo-livro}.

% Citação DIRETA LONGA (mais de 3 linhas): destacada em parágrafo próprio, com
% recuo de 4 cm da margem esquerda, fonte menor, espaço simples e SEM aspas
% (NBR 10520:2023; guia UNIVALI). Use o ambiente \begin{citacao}...\end{citacao}:
% o abnTeX2 já aplica recuo (4 cm), \ABNTEXfontereduzida e espaçamento simples.
% A chamada (autor, ano, página) precede o bloco; não se usam aspas.
Já a citação com mais de três linhas é destacada do texto, como em
\textcite[p. 86]{exemplo-livro}:

\begin{citacao}
  Lorem ipsum dolor sit amet, consectetur adipiscing elit, sed do eiusmod tempor
  incididunt ut labore et dolore magna aliqua. Ut enim ad minim veniam, quis
  nostrud exercitation ullamco laboris nisi ut aliquip ex ea commodo consequat.
  Duis aute irure dolor in reprehenderit in voluptate velit esse cillum dolore eu
  fugiat nulla pariatur.
\end{citacao}

% Variações de chamada (NBR 10520:2023): grifo nosso (ao destacar parte da
% citação), supressão com reticências entre colchetes [...] (trecho omitido) e
% apud / citação de citação ("Autor original (ano apud Autor consultado, ano,
% p.)"). A nota de rodapé vai em fonte menor e espaço simples (NBR 14724:2024).
Ao destacar parte de uma citação, indica-se o grifo na
chamada~\parencite[p.~22, grifo nosso]{exemplo-artigo}; trechos omitidos são
marcados com reticências entre colchetes (\enquote{[\ldots] texto restante}).
Quando se cita uma fonte por meio de outra, usa-se \emph{apud}: para Autor
Original (1998 apud \citeauthor{exemplo-dissertacao}, \citeyear{exemplo-dissertacao},
p.~30), o conceito se mantém.\footnote{Esta é uma nota de rodapé: digitada dentro
das margens, em fonte menor e espaço simples, separada do texto por um filete
(NBR 14724:2024).}

% --- EXEMPLO DE FIGURA -----------------------------------------------
% Para usar uma imagem real, troque o conteúdo do \centering pela linha:
%   \includegraphics[width=0.75\textwidth]{imagens/sua-figura.png}
% O comando \figuraplaceholder (definido em tex/_placeholders.tex) apenas
% desenha uma caixa cinza para o template compilar sem imagens externas.
\begin{figure}[!htb]
  \caption{Legenda da figura (acima da figura, conforme NBR 14724).}
  \label{fig:exemplo}
  \centering
  % A fonte fica alinhada à margem esquerda DA FIGURA (não do texto).
  \elementocomfonte{\figuraplaceholder{0.6\textwidth}{4cm}{imagens/sua-figura.png}}%
    {Fonte: Elaborada pelo autor (\the\year{}).}
\end{figure}

%-------------------------------------------------------------
\section{Formulação matemática}
\label{sec:formulacao}
%-------------------------------------------------------------

Equações são numeradas e podem ser referenciadas, como mostra a
\autoref{eq:exemplo}:

\begin{equation}
  E = m c^{2},
  \label{eq:exemplo}
\end{equation}

\noindent onde $E$ é a energia, $m$ a massa e $c$ a velocidade da luz no vácuo.
At vero eos et accusamus et iusto odio dignissimos ducimus.

%-------------------------------------------------------------
\section{Síntese}
%-------------------------------------------------------------

Os principais pontos abordados foram:

\begin{enumerate}[label=(\roman*)]
  \item primeiro ponto: lorem ipsum dolor sit amet;
  \item segundo ponto: consectetur adipiscing elit;
  \item terceiro ponto: sed do eiusmod tempor incididunt.
\end{enumerate}

% Alíneas (ABNT NBR 6024): subdivisão SEM título dentro de uma seção, quando não
% se quer numerar com \section. Introduzidas por frase terminada em dois-pontos;
% cada alínea em letra minúscula, terminando em ";" (a última em "."), ordenadas
% a) b) c). Ambiente \begin{alineas} (abnTeX2); subdivisão extra: \begin{subalineas}.
Sem necessidade de títulos, a enumeração também pode ser feita por alíneas:

\begin{alineas}
  \item primeira alínea, em letra minúscula, terminando em ponto-e-vírgula;
  \item segunda alínea, que pode se desdobrar em subalíneas:
  \begin{subalineas}
    \item primeira subalínea;
    \item segunda subalínea;
  \end{subalineas}
  \item última alínea, terminando em ponto.
\end{alineas}
